Since I need a short and effective method of checking the files in a remote environment without tethered access (USB connection), I opted for saving the files on an uSD card. Although this does not provide real-time imagery, the slightly smaller scope allows us to better focus on the functionality of the camera rather than how the fully processed image is moved around.
\nl
Instead of PCB-mounting an uSD card directly, I chose the \href{https://au.mouser.com/ProductDetail/Adafruit/254?qs=GURawfaeGuAkwqCF4BmPzA%3D%3D&srsltid=AfmBOoqN0VaaHsFOYARHTvz8zYqY9TRrMM7HXF_dcSo8h28U_mns63k3}{254 Adafruit} breakout board because it can be easily connected via standard PCB headers. Although it appears this step has been shortcut, the prototype will not use an uSD card reader so it suffices as an effective solution without having to deal with the techincal aspects.
As the prototype is expected to offload images to another device, this breakout-based approach approach works well for viewing the file in its intermediate stage (leaving the processor and being sent to FFS).
